\chapter*{Resumen}
\addcontentsline{toc}{chapter}{Resumen}

El email marketing es uno de los canales publicitarios digitales con mayor retorno potencial, pero su eficacia depende de la capacidad de dirigirse al usuario adecuado con el mensaje correcto. Las empresas que operan en este canal se enfrentan a un problema de asignación de inversión: ¿qué usuarios tienen mayor probabilidad de interactuar con un producto determinado, y cómo evoluciona esa demanda a lo largo del tiempo?

Este trabajo aborda ese problema mediante un \textit{framework} integral de aprendizaje automático aplicado a campañas reales del mercado español. El pipeline desarrollado cubre desde la ingesta y limpieza de datos hasta la generación de recomendaciones personalizadas, pasando por el enriquecimiento demográfico y geográfico de perfiles de usuario, la segmentación no supervisada de clientes y el modelado de propensión.

El sistema se estructura en cuatro módulos. El primero (M0/M1) realiza dos niveles de segmentación: una segmentación demográfica pura mediante autoencoder y KMeans orientada al arranque en frío, y una segmentación comportamental sobre el historial de interacciones mediante PCA y KMeans. El segundo módulo (M2) desarrolla un modelo de propensión basado en LightGBM que incorpora variables de usuario y de producto ---incluidos \textit{embeddings} del texto del asunto---, con corrección explícita del sesgo de muestreo; se comparan cinco algoritmos (incluida una red neuronal) y dos arquitecturas (un modelo plano frente a uno jerárquico sector\,$\rightarrow$\,producto), y se adopta el modelo plano, que supera al jerárquico en todos los algoritmos. El tercer módulo (M3) implementa un recomendador por cluster: sirve a cada usuario el \textit{ranking} de categorías de su segmento comportamental. El cuarto módulo (M4) aplica el modelo Prophet para la predicción del volumen semanal de clics.

La evaluación se realiza sin fuga de datos (validación agrupada por usuario en M2; partición temporal y reconstrucción del segmento solo con \textit{train} en M3) y con contraste estadístico (intervalos \textit{bootstrap}, Friedman, Wilcoxon y McNemar con corrección de Holm). Un hallazgo central es que la señal útil para la propensión reside en el contenido del producto (PR-AUC, área bajo la curva de precisión-exhaustividad, 0,402 del modelo con variables de producto frente a 0,313 del baseline), mientras que las variables demográficas ---sintéticas--- tienen un poder predictivo limitado. En recomendación, una vez eliminada la fuga, la popularidad es un \textit{baseline} que el segmento no logra batir (Hit~Rate@5 0,586 del recomendador frente a 0,629 de la popularidad); el valor del recomendador por cluster es operativo (interpretabilidad y arranque en frío), no de mejora del \textit{ranking}.

\bigskip
\noindent\textbf{Palabras clave:} email marketing, modelado de propensión, LightGBM, segmentación de usuarios, autoencoder, sistemas de recomendación, enriquecimiento demográfico, corrección de prior shift.

\bigskip\bigskip

\noindent\textbf{Abstract}

\medskip
Email marketing is one of the digital advertising channels with the highest potential return on investment, but its effectiveness depends on the ability to target the right user with the right message. Companies operating in this channel face an investment allocation problem: which users are most likely to engage with a given product, and how does that demand evolve over time?

This thesis addresses that problem through a comprehensive machine learning framework applied to real campaigns from the Spanish market. The developed pipeline covers the full lifecycle from data ingestion and cleaning to personalised recommendations, including demographic and geographic enrichment of user profiles, unsupervised customer segmentation, and propensity modelling.

The system is organised into four modules. The first (M0/M1) performs two levels of segmentation: a purely demographic segmentation via autoencoder and KMeans for cold-start assignment, and a behavioural segmentation over interaction history via PCA and KMeans. The second module (M2) develops a LightGBM propensity model that incorporates both user and product features ---including \textit{embeddings} of the email subject---, with explicit correction of the sampling bias; five algorithms (including a neural network) and two architectures (a flat model vs.\ a hierarchical sector\,$\rightarrow$\,product one) are compared, and the flat model ---which outperforms the hierarchical one across all algorithms--- is adopted. The third module (M3) implements a cluster-level recommender: each user is served the category \textit{ranking} of their behavioural segment. The fourth module (M4) applies the Prophet model for weekly click-volume forecasting. Evaluation is performed without data leakage (user-grouped cross-validation in M2; temporal split and train-only reconstruction of the segment in M3) and with statistical testing (\textit{bootstrap} intervals, Friedman, Wilcoxon and McNemar with Holm correction). A central finding is that the useful signal for propensity lies in product content (PR-AUC 0.402 for the product-aware model vs.\ 0.313 for the baseline), whereas the ---synthetic--- demographic features have limited predictive power. In recommendation, once the leakage is removed, popularity is a baseline the segment fails to beat (Hit~Rate@5 0.586 for the recommender vs.\ 0.629 for popularity); the value of the cluster-level recommender is operational (interpretability and cold-start), not ranking improvement.

\bigskip
\noindent\textbf{Keywords:} email marketing, propensity modelling, LightGBM, user segmentation, autoencoder, recommender systems, demographic enrichment, prior shift correction.
